\documentclass[../main.tex]{subfiles}

\begin{document}

\section{Chapter 9: Storytelling and Presentation Narrative}
\label{sec:data-storytelling-design-opening}

This section analyzes the Vietnam Financial Inclusion Dashboard 2025 as a data
story rather than as a collection of independent Power BI visuals. Since the
dataset and sampling background have already been discussed in the previous
sections, this chapter focuses on how the dashboard transforms weighted survey
evidence into a policy-oriented narrative about financial inclusion, digital
readiness, account access, cash-flow usage, and financial resilience in
Vietnam.

\dashboardxref{fig:dashboard-overview}

\subsection{Dashboard Purpose and Storytelling Lens}
\label{subsec:dashboard-purpose-storytelling-lens}

Within the established data context, the dashboard functions as an institutional
policy workspace. It does not only report whether adults have financial
accounts; it also asks whether people are digitally ready, whether account
ownership becomes actual financial usage, and whether financial participation is
connected to saving, borrowing, and resilience outcomes. For this reason, the
dashboard is best read as a sequential policy story rather than as a set of
disconnected descriptive charts.

The organizing logic of the dashboard can be summarized as:

\begin{center}
\textbf{Context \(\rightarrow\) Readiness \(\rightarrow\) Access \(\rightarrow\) Usage \(\rightarrow\) Resilience}
\end{center}

This structure is important because financial inclusion is not a single-step
condition. Digital access does not always lead to active digital verification;
account ownership does not always lead to digital cash-flow usage; and
financial participation does not automatically guarantee resilience during
financial shocks. The dashboard therefore guides the reader through a
multi-stage diagnosis.

\storyroadmapfigure{fig:story-roadmap}

\subsection{Narrative Walkthrough of the Dashboard}
\label{subsec:narrative-walkthrough}

Before examining the formal storytelling dimensions, this walkthrough briefly
shows how each dashboard page contributes to the overall story.

\subsubsection{Page 1: Establishing the Socioeconomic Context}
\label{subsubsec:page1-context}

Page 1 acts as the exposition of the story by defining who is being analyzed
before introducing financial indicators. The KPI cards show an Urban
Concentration Rate of 30.33\%, a Secondary Education Threshold Rate of 89.98\%,
and an In Workforce Rate of 87.92\%. These values establish the demographic and
socioeconomic baseline for the later financial inclusion analysis.

The strongest visual evidence on this page is the socioeconomic matrix. The
Poorest 20\% group has only 6 respondents with tertiary education or more,
while the Richest 20\% group has 72. This contrast frames financial inclusion
as a socially structured issue shaped by income, education, location, and
workforce participation.

\dashboardxref{fig:socioeconomic-matrix}

\subsubsection{Page 2: Diagnosing Digital Readiness}
\label{subsubsec:page2-readiness}

Page 2 shifts the story from population context to digital readiness. Internet
Use in the last three months is 87.85\%, and Internet Use in the last seven
days is 85.39\%, suggesting that internet access is frequent rather than
occasional. However, the central insight is the gap between digital access and
digital activation: Foundational ID coverage reaches 99.65\%, Online Digital ID
reaches 71.19\%, but Active Digital ID Verification is only 32.68\%.

This page therefore shows that digital infrastructure alone is not enough.
Although many adults are connected and identifiable, far fewer actively use
digital verification. This gap becomes an important bridge between basic
readiness and meaningful digital financial participation.

\begin{figure}[H]
    \centering
    \begin{subfigure}[t]{0.31\textwidth}
        \centering
        \includegraphics[width=\textwidth]{fig-page2-connectivity-depth.png}
        \caption{Connectivity depth: 3M vs. 7D.}
        \label{fig:page2-connectivity-depth}
    \end{subfigure}
    \hfill
    \begin{subfigure}[t]{0.31\textwidth}
        \centering
        \includegraphics[width=\textwidth]{fig-page2-digital-vs-foundational-id.png}
        \caption{Digital vs. foundational ID coverage.}
        \label{fig:page2-id-coverage}
    \end{subfigure}
    \hfill
    \begin{subfigure}[t]{0.31\textwidth}
        \centering
        \includegraphics[width=\textwidth]{fig-page2-active-digital-id.png}
        \caption{Active digital ID verification rate.}
        \label{fig:page2-active-digital-id}
    \end{subfigure}
    \caption{Page 2 readiness evidence used in the storytelling analysis.}
    \label{fig:page2-readiness-overview}
\end{figure}

\subsubsection{Page 3: Identifying Account Access and Exclusion}
\label{subsubsec:page3-access}

Page 3 is the central diagnostic layer of the dashboard. It moves from digital
readiness to direct financial access. The KPI cards report Banking Inclusion at
74.93\%, Financial Institution Account at 74.05\%, Mobile Money Account at
46.69\%, and Digital Account at 67.03\%. The 27.36 percentage-point gap between
financial institution accounts and mobile money accounts suggests that
traditional accounts remain more widespread than mobile money.

The page also turns access into an exclusion diagnosis. Its filters allow users
to segment the story by Urbanicity, Gender, Workforce Status, Education Level,
and Income Group. Meanwhile, the Exclusion Friction Log identifies barriers such
as not enough money, distance, fees, lack of trust, and lack of documentation.
As a result, Page 3 does not only ask how many people are included; it also
asks who remains excluded and why.

\dashboardfigure{0.82}{fig-page3-income-gender.png}
{Banking inclusion by income group and gender.}
{fig:page3-income-gender}

\noindent\textit{See also Figures~\ref{fig:account-kpi-cards},
\ref{fig:page3-filter-bar}, \ref{fig:account-composition-flow}, and
\ref{fig:exclusion-friction-log}.}

\subsubsection{Page 4: Moving from Access to Cash-Flow Usage}
\label{subsubsec:page4-usage}

Page 4 prevents the story from stopping at account ownership. It asks whether
financial flows are actually delivered through digital channels or still depend
on cash. Corporate/Formal Wages are largely digital, with 80.99\% delivered
directly to an account and 18.60\% cash-only. In contrast, Agricultural Sales
remain cash-dependent, with 73.42\% cash-only and only 25.79\% delivered
digitally to an account.

The page also reveals a geographic usage gap. Urban Merchant Payments reach
65.79\%, compared with 50.30\% in rural areas; Urban In-Store Digital Pay
reaches 62.67\%, compared with 46.11\% in rural areas. These comparisons show
that financial inclusion must be evaluated through actual transaction behavior,
not only through account ownership.

\dashboardfigure{0.56}{fig-page4-digital-merchant-payments.png}
{Digital merchant and in-store payment by geography.}
{fig:page4-digital-merchant-payments}

\dashboardxref{fig:national-inflows}

\subsubsection{Page 5: Evaluating Financial Health and Resilience}
\label{subsubsec:page5-resilience}

The final page connects financial participation to financial health and
resilience. The KPI cards show that 73.05\% of respondents saved, while 45.98\%
borrowed. This shifts the dashboard from access and usage toward financial
outcomes.

The resilience tables provide the strongest evidence for this final stage. For
the Poorest 20\% group, Family \& Friends is 86 while Formal Savings is only
30. For the Richest 20\% group, Formal Savings is 60 while Family \& Friends is
34. This contrast suggests that poorer groups depend more heavily on informal
support, while richer groups are more connected to formal savings channels.

\dashboardfigure{0.92}{fig-page5-resilience-overview.png}
{Financial Health \& Resilience page overview.}
{fig:page5-resilience-overview}

\dashboardfigure{0.68}{fig-page5-coping-income-group.png}
{Coping channels by income group.}
{fig:page5-coping-income-group}

\subsection{Bridge to the Storytelling Analysis}
\label{subsec:bridge-to-storytelling-analysis}

Together, the five pages form a clear policy narrative. Page 1 establishes the
population context, Page 2 tests digital readiness, Page 3 diagnoses account
access and exclusion, Page 4 examines cash-flow usage, and Page 5 evaluates
financial health and resilience. This order prevents financial inclusion from
being reduced to a single account ownership metric.

The remainder of this chapter analyzes the dashboard through four storytelling
dimensions:

\begin{enumerate}
\item \textbf{Story structure}: how the dashboard organizes the narrative from context to resilience.
\item \textbf{Storytelling principles}: how overview-first design, segmentation, comparison, gap analysis, and diagnostic framing support the story.
\item \textbf{Narrative style}: how the dashboard uses a neutral, institutional, evidence-based policy tone.
\item \textbf{Interaction and exploration}: how filters and subgroup views allow users to move from national indicators to targeted policy questions.
\end{enumerate}

\dashboardxref{fig:story-roadmap}

\end{document}
